\section{Counting comparisons and paying for moves}
\label{sec:accounting}

We analyze a completed play with $q$ comparisons. If $q>2T$ the desired
bound already holds, so assume $q\leq2T$.
A cost is a decrease of $\logtwo Z$; an admission gain has negative cost.
We keep these signs throughout. Charging a fixed positive amount for each
move would lose the factorial cancellation that makes the argument work.

\subsection{Population, incidence, and separation}

For a node $v$, let $s_v$ be the number of items ever entering it, and
$k_{vc}$ the number entering its child at $c$. Set
\[
 p_v=\sum_c k_{vc},\qquad r_v=s_v-p_v,\qquad
 \Mov=\sum_v p_v .
\]
Here $r_v$ counts the items that end free in $v$, and $\Mov$ is the total
number of moves.

\begin{lemma}[Population and incidence]
\label{lem:population}
There are at most $2n-1$ nonempty nodes, and
\begin{equation}
 \sum_v r_v=n,\qquad \sum_v s_v=n+\Mov,\qquad \Mov\leq2q.
\label{eq:population}
\end{equation}
At most $K_v=2s_v+2$ positions are ever created in node $v$, and hence
\begin{equation}
                    \Ksum:=\sum_vK_v\leq6n+2\Mov.
\label{eq:palettes}
\end{equation}
At most $q$ active edges are retired.
\end{lemma}
\begin{proof}
Each item ends free at exactly one node and contributes one new node
occurrence per move. This proves the first two identities.
Every moved item has an incident comparison edge, either at movement or
when a path containing it was reserved. Charge that move to the corresponding
edge endpoint. An item leaves a node only once, so an edge receives at most
two charges. Each edge came from an actual comparison at that node.
This gives $\Mov\leq2q$, bounds edge retirements, and also rules out an
infinite repair sequence during a finite transcript.

The two endpoints of any charged edge cannot both enter the same child:
the second would be publicly comparable with an old child item.
A node with only one child therefore retains an item. So does a leaf.
Those retained items are distinct across nodes. A finite rooted tree has
fewer branching nodes than leaves, giving at most $2n-1$ nodes.
In particular each child is a proper subset of its parent.

Initial positions and top-position additions create at most $s_v+2$
positions, and at most one endpoint insertion occurs per moved item.
Thus $K_v=2s_v+2$ is an upper bound. Summing proves \eqref{eq:palettes}.
\end{proof}

Define the separation budget
\begin{equation}
                     E_{\mathrm{sep}}
                     =\sum_{v,c}k_{vc}\logtwo\frac{s_v}{k_{vc}}.
\label{eq:sep}
\end{equation}
This quantity is used only in the analysis.

\begin{lemma}[One separation budget]
\label{lem:entropy}
\[
                   E_{\mathrm{sep}}
                   =T-\sum_v r_v\logtwo s_v\leq T.
\]
\end{lemma}
\begin{proof}
The contribution of $v$ is
$s_v\logtwo s_v-\sum_c k_{vc}\logtwo k_{vc}-r_v\logtwo s_v$.
Every nonroot $s_v$ equals the multiplicity on its incoming edge.
The first two terms telescope, leaving only the root term $T$.
\end{proof}

\begin{lemma}[Expensive moves are rare]
\label{lem:rare}
Let $\Ord,\Sc$ count ordinary and scale moves. Then
\begin{equation}
 \Ord\leq\frac{2q}{D-1}+\frac{T}{\logtwo H},
 \qquad
 \Sc\leq\frac{T}{\logtwo(J/2)}.
\label{eq:rare}
\end{equation}
\end{lemma}
\begin{proof}
A degree move has at least $D-1$ incident edges when selected. Charging
all those endpoints and using the same incidence argument bounds their
number by $2q/(D-1)$.

For destination counts $k_c$ totaling $p$, partition the admissions into
groups. In a group of size $g$, let $t_c$ go to destination $c$.
Concavity of entropy gives
\begin{equation}
            \sum_{\mathrm{groups}}\sum_c t_c\logtwo\frac{g}{t_c}
            \leq\sum_c k_c\logtwo\frac{p}{k_c}.
\label{eq:group-entropy}
\end{equation}
This follows by applying concavity of $-x\logtwo x$ to the average of the
groups' destination-frequency vectors, weighted by $g/p$.

All $H$ items of a reserved path eventually move, even if some do so during
width repair. They are publicly comparable and therefore enter distinct
children. Each such group contributes $H\logtwo H$ to the left side of
\eqref{eq:group-entropy}. The groups are disjoint at a node. The right
side is bounded by its contribution to \eqref{eq:sep}, since $p_v\leq s_v$.
Lemma~\ref{lem:entropy} bounds the number of reserved path members by
$T/\logtwo H$, and hence also bounds their ordinary moves.

Apply the same argument separately to scale groups. Each has $J$ items
and destination multiplicities at most two, so contributes at least
$J\logtwo(J/2)$. This proves the other bound.
\end{proof}

\subsection{Representative reuse}

\begin{lemma}[A lifetime dictionary of representatives]
\label{lem:dictionary}
Suppose at most $K$ positions are ever born in a node. From an interval of
$m$ positions each having probability at least $1/(Vm)$, the representative
rule choosing block size $b$ retains at least the fraction
\eqref{eq:representative-cost}, with an absolute constant $C$.
Over the node's lifetime, the number of different destinations selected
with size $b$ is at most
\begin{equation}
                       \frac Kb\bigl(2+\logtwo(K/b)\bigr).
\label{eq:dictionary}
\end{equation}
The interval-size bound may be taken as
$m\leq8b(2+\logtwo(K/b))$, including the restricted scale-group choices.
\end{lemma}
\begin{proof}
A consecutive subsequence of a batch has at most two maximal aligned blocks
of any one size: each such block has a parent crossing one of the two
endpoints. Thus blocks of size at most $b$ cover fewer than $4b$ members
of any batch. Batches smaller than $b$ have fewer than $b$ members in total,
and there are at most $1+\logtwo(K/b)$ other current batches.
Endpoint trimming retains a constant fraction of $m$, giving the stated
bound. Intervals of size less than eight are covered by $b=1$.
The exclusions in a scale group remove less than a quarter of the central
interval by the width ratios in Lemma~\ref{lem:width-repairs}; the same
constant eight suffices.

At a fixed batch size $B'$, at most $K/B'$ batches ever occur, since their
birth sets are disjoint. Each has $B'/b$ aligned blocks of size $b$.
Summing over $B'=b,2b,\ldots$ gives \eqref{eq:dictionary}.
One additional $K$ term at size one covers exceptional endpoint choices.
\end{proof}

\subsection{Cheap moves and factorial credits}

Credits can be separated into cheap, ordinary, and scale admissions.
At the $j$th admission of one category, the child's total admission number
is at least $j$. Thus $k$ admissions of that category supply a factor
at least $k!$, by Lemma~\ref{lem:birth}.

For $k$ cheap moves to one child, the net cost is at most
\[
      k\logtwo G-\logtwo(k!)=\logtwo\frac{G^k}{k!}
                          \leq G\logtwo e,
\]
because $G^k/k!$ is a term of $\exp(G)$.
There are at most $2n-2$ children, and $2\logtwo e<3$. Therefore
\begin{equation}
                     \text{net cheap cost}\leq3nG.
\label{eq:cheap}
\end{equation}
Cheap moves need not be rare.

\subsection{The signed cost of the other moves}
\label{sec:signed}

Put $X=\Ord+\Sc$ and $\sigma=1+2B$, so $2\cdot4^B=2^\sigma$.
Set aside a surcharge $\sigma$ for every scale move. Lemma~\ref{lem:dictionary}
then bounds the remaining price of either kind of move by
$\logtwo[128b(2+\logtwo(K_v/b))]$. An ordinary move uses $V=16$;
the scale surcharge covers its larger $V$.

\begin{lemma}[Total net move cost]
\label{lem:net}
All fixing costs minus all admission gains total at most
\begin{equation}
 3nG+\sigma\Sc+
 40X\logtwo\left(2+\frac{n+T+\Mov}{X}\right).
\label{eq:net}
\end{equation}
The last term is defined as zero if $X=0$. The constant $40$ is independent
of every stage parameter.
\end{lemma}
\begin{proof}
Fix a node and one of the two expensive categories. For each destination
$c$, let $b_c$ be its largest selected block size in this category, and
$j_c=\logtwo(K_v/b_c)$. By \eqref{eq:dictionary}, the positive weights
\[
                 Q_c=\frac{b_c}{K_v(2+j_c)^3}
                 \quad\hbox{satisfy}\quad \sum_c Q_c\leq1.
\]
Indeed, destinations with largest block size $b$ contribute at most
$(2+\logtwo(K_v/b))^{-2}$, and these arguments form a unit-spaced grid
starting at a nonnegative number. Their sum is less than
$\sum_{r\geq0}((r+1)(r+2))^{-1}=1$.

Let $k_c$ be this category's multiplicity at $c$.
The price increases with block size. Using
$k!\geq(k/e)^k$ and $e<4$, its net cost after the factorial credit is at most
\begin{equation}
                       k_c\logtwo\frac{512b_c(2+j_c)}{k_c}.
\label{eq:destination-price}
\end{equation}
Discard negative terms and let $I$ index the remaining
node--destination--category triples, with $Y=\sum_I k_c\leq X$.
If $Y=0$ there is nothing to bound. Positivity of
\eqref{eq:destination-price} gives
\[
 k_c\leq512K_v2^{-j_c}(2+j_c),\qquad
 j_c\leq2\logtwo(K_v/k_c)+22,
\]
where the second estimate uses $\logtwo(2+j)\leq j/2+2$.

For positive $a_i,k_i$, concavity of the logarithm gives
\begin{equation}
             \sum_i k_i\logtwo(a_i/k_i)
             \leq Y\logtwo\left(\frac{\sum_i a_i}{Y}\right),
             \qquad Y=\sum_i k_i .
\label{eq:logsum}
\end{equation}
Since $b_c(2+j_c)=K_vQ_c(2+j_c)^4$ and there are two categories,
summing \eqref{eq:destination-price} yields
\begin{equation}
  \text{remaining cost}\leq
  Y\logtwo\frac{1024\Ksum}{Y}+4\sum_I k_c\logtwo(2+j_c).
\label{eq:price-sum}
\end{equation}

Write $k^*_{vc}$ for total admissions to that child, including all categories.
Set $U=n+T+\Mov$. Since $K_v\leq4s_v$,
\begin{align*}
\sum_I k_c\logtwo(K_v/k_c)
&=\sum_I k_c\logtwo(K_v/s_v)
 +\sum_I k_c\logtwo(s_v/k^*_{vc})
 +\sum_I k_c\logtwo(k^*_{vc}/k_c)\\
&\leq2Y+T+2\Mov.
\end{align*}
The middle sum is bounded by Lemma~\ref{lem:entropy}, since the selected
category multiplicities sum to at most $k^*_{vc}$.
For the last sum use $x\logtwo(1/x)\leq1$ on each of the two categories.
By \eqref{eq:logsum} and the bound on $j_c$,
\[
\begin{split}
\sum_I k_c\logtwo(2+j_c)
&\leq Y\logtwo\left(24+\frac2Y\sum_I k_c\logtwo(K_v/k_c)\right)\\
&\leq Y\logtwo(28+4U/Y)\\
&\leq5Y+Y\logtwo(2+U/Y).
\end{split}
\]
Also $\Ksum\leq6U$. Substitute these estimates in
\eqref{eq:price-sum} to obtain
\[
 \text{remaining cost}\leq33Y+5Y\logtwo(2+U/Y)
                      \leq40Y\logtwo(2+U/Y).
\]
The function $y\logtwo(2+U/y)$ increases for $y>0$ and tends to zero
as $y\downarrow0$, so replace $Y$ by $X$.
Adding \eqref{eq:cheap} and $\sigma\Sc$ proves the claim.
\end{proof}

The weights $Q_c$, maximum block sizes, and final multiplicities are proof
devices. The online algorithm uses only the currently available blocks.

\subsection{The terminal count and the comparison inequality}

\begin{lemma}[Terminal count]
\label{lem:terminal}
If the sorter has finished and $n\geq2$, $q\leq2T$, then
\[
                      \logtwo Z_{\mathrm{final}}
                      \leq n\logtwo(21+24\logtwo n).
\]
\end{lemma}
\begin{proof}
The public order is total. The $r_v$ free items at $v$ must therefore
occupy a nondecreasing sequence of positions in their known order.
Ignoring additional bounds gives at most $\binom{K_v+r_v-1}{r_v}$
assignments. Concatenate the node position lists in a fixed order.
The tuple of node sequences injects into a nondecreasing sequence of
$n=\sum_v r_v$ positions from the concatenated list: its pieces are
recoverable from the known $r_v$.
Consequently
\[
 Z_{\mathrm{final}}\leq\binom{\Ksum+n-1}{n}
              \leq\left(\frac{e(\Ksum+n)}n\right)^n .
\]
Use $\Ksum\leq6n+2\Mov$, $\Mov\leq2q\leq4T$, and $e<3$.
\end{proof}

\begin{theorem}[Fixed-stage comparison bound]
\label{thm:fixed}
Every completed play with $n\geq2$, $q\leq2T$ satisfies
\begin{equation}
\begin{split}
T\leq{}&(1+15\eta)q+n\logtwo(21+24\logtwo n)+3nG+\sigma\Sc\\
       &\quad+40X\logtwo\left(2+\frac{n+T+\Mov}{X}\right),
\end{split}
\label{eq:master}
\end{equation}
where $\Mov\leq2q$ and \eqref{eq:rare} bounds $X=\Ord+\Sc$.
\end{theorem}
\begin{proof}
Each actual comparison retains either all assignments or at least
$1/2-2\eta$ at one node. There are at most $q$ strict edge cuts, each
retaining at least $1-\eta$. For $\eta\leq1/16$, their combined cost is
at most $(1+15\eta)q$, using $-\ln(1-z)\leq2z$ for $z\leq1/2$.
Initially $\logtwo Z=T$. Endpoint insertions preserve old assignments,
and Lemma~\ref{lem:net} includes the signed admission credits.
Telescope all changes and apply Lemma~\ref{lem:terminal}.
\end{proof}
